% Въ лѣто Господне ҂вк҃ѕ.
%
% Господи Іисусе Христе, Сыне Божій, помилуй меня грѣшнаго.

\documentclass[17pt]{extarticle}

\usepackage[a5paper, portrait, margin=.7cm, top=0cm, bottom=0cm]{geometry}

%\usepackage{pgfpages}
%\pgfpagesuselayout{2 on 1}[a4paper, landscape]

\usepackage{fontspec}
\usepackage{xkeyval}
\usepackage{microtype}

\usepackage{polyglossia}
\setmainlanguage{churchslavonic}

\usepackage{churchslavonic}
\usepackage{lettrine}
\usepackage{graphicx}  % for \scalebox

\usepackage{hyperref}
\hypersetup{
	pdftitle={Икона Божией Матери Толгская, молитва},
	pdfauthor={Православная Церковь}
}


\makeatletter

\def\@texttop{\vfill}
\def\@textbottom{\vfill}

\makeatother


\newcommand{\cuL}[1]{\scalebox{1.5}{\cuKinovarColor#1}}
\newcommand{\cul}[1]{\mbox{\cuKinovarColor#1}}

\newfontfamily\noto{Noto Sans Symbols 2}
\newfontfamily\churchslavonicfont[Script=Cyrillic,Ligatures=TeX,HyphenChar=_]{Triodion}
%\newfontfamily\cyrillicfontsf{DejaVu Sans}

\hyphenpenalty=10000
\setlength\emergencystretch{4em}
\def\hyph{{\_}{}{}\penalty-10000}

%\linespread{1.2}
\setlength{\parindent}{1cm}
%\setlength{\parskip}{8pt}

\newcommand{\sub}[1]{\vspace{4pt}{\centering\cuKinovar{#1}\par}}


\begin{document}

\pagestyle{empty}

\begin{center}
\cuKinovarColor{\Large\noto 🙚}
\begin{minipage}{.57\textwidth}\centering
	\vspace{-6pt}
	Прес҃тѣ̀й бцⷣе пред̾ і҆кѡ́ною є҆ѧ́, \\\vspace{-4pt} именꙋ́емыѧ то́лгскаѧ, моли́тва:
\end{minipage}
{\Large\noto 🙘}
\end{center}

\vspace{-8pt}
\cuL Прїими́, ѽ всебл҃гослове́ннаѧ и҆ всемо́щнаѧ пречⷭ҇таѧ гжⷭ҇ѣ̀ влⷣчце, бцⷣѣ̀ дв҃о, сїѧ̑ моли̑твы, со слеза́ми тебѣ́ ны́нѣ приноси́мыѧ ѿ на́съ, недосто́йныхъ ра̑бъ твои́хъ, къ цѣльбоно́сномꙋ ѻ҆́бразꙋ твоемꙋ̀ со ѹ҆миле́нїемъ притека́ющихъ, ꙗ҆́кѡ тебѣ́ само́й тꙋ̀ сꙋ́щей, и҆ слы́шащей моле́нїѧ на̑ша, и҆ подаю́щей съ вѣ́рою просѧ́щымъ по кое́мꙋждо проше́нїю и҆сполне́нїе. Скорбѧ́щымъ бо ско́рби ѡ҆блегча́еши, немощны̑мъ здра́вїе да́рꙋеши, раз\hyph сла́бленныхъ и҆ недꙋ́жныхъ и҆сцѣлѧ́еши, ѿ бѣ́сныхъ бѣ́сы прогонѧ́еши, ѡ҆би́димыѧ ѿ ѡ҆би́дъ и҆збавлѧ́еши, и҆ наси́лꙋємыѧ спаса́еши, грѣ́шныхъ ка́ющихсѧ проща́еши, прокаже́нныхъ ѡ҆чища́еши, и҆ ма́лыхъ дѣте́й ми́лꙋеши, ѿ ѹ҆́зъ и҆ темни́цъ свобожда́еши, и҆ всѧ́кїѧ многоразли́чныѧ стра́сти врачꙋ́еши ты́, гжⷭ҇ѣ̀ влⷣчце бцⷣе: всѧ̑ бо сꙋ́ть возмѡ́жна хода́тайствомъ твои́мъ къ сн҃ꙋ твоемꙋ̀, хрⷭ҇тꙋ̀ бг҃ꙋ на́шемꙋ.

\vspace{1pt}
\cuL Ѽ всепѣ́таѧ мт҃и, прест҃а́ѧ бцⷣе дв҃о мр҃і́е! не преста́й моли́тисѧ ѡ҆ на́съ, недосто́йныхъ рабѣ́хъ твои́хъ, сла́вѧщихъ тѧ̀ и҆ почита́ющихъ, и҆ покланѧ́ющихсѧ пречⷭ҇томꙋ ѻ҆́бразꙋ твоемꙋ̀, и҆ наде́ждꙋ и҆мꙋ́щихъ невоз\hyph вра́тнꙋ и҆ вѣ́рꙋ несꙋмнѣ́ннꙋ къ тебѣ́, приснодв҃ѣ пре\hyph сла́внѣй и҆ непоро́чнѣй, ны́нѣ и҆ прⷭ҇нѡ и҆ во вѣ́ки вѣкѡ́въ.

\vspace{-6pt}
\cuL{А҆}ми́нь.
\vspace{6pt}

\end{document}
